% templates/conclusions_and_diagnostics.tex.mustache
% AUTO-GENERATED Pipeline Verification and Analytical Conclusions
% Rendered by Mustache template from curated diagnostics
% DO NOT EDIT MANUALLY

\section{System Diagnostics and Definitive Conclusions}

\subsection{Pipeline Execution Verification}

The execution of the \texttt{SpectralBL-Analytics} pipeline across divergent data scales demonstrates the mathematical robustness of the low-rank attractor manifold projection. The framework successfully achieves structural decomposition without reliance on local gradient Richardson number ($\text{Ri}_g$) formulations, which typically fail under stable conditions.

\subsection{Comparative Information-Theoretic Summary}

Based on the processed analytics run, the structural characteristics of the underlying dataset are summarized below:

\begin{table}[htbp]
    \centering
    \small
    \begin{tabular}{l p{9cm}}
        \hline
                \textbf{Diagnostic Metric} & \textbf{Pipeline Evaluation Summary for \texttt{ FLOSS }} \\ \hline
        Data Footprint & Analyzed 70796 observations across a temporal scale of 11732.5 hours. \\
        Manifold Complexity & Mean singular value entropy registers at $H_{\text{mean}} = 0.425$, yielding a Shannon effective dimension of $D_{\text{eff}} = 1.5369$. \\
        Numerical Stability & The condition number $\kappa = 48.95$ confirms a well-conditioned low-rank basis projection ($N_0 = 0$ nullspace modes). \\ \hline
    \end{tabular}
    \caption{Key pipeline execution and manifold metrics for the current campaign.}
    \label{tab:pipeline_summary_FLOSS}
\end{table}

\subsection{Definitive Scientific Conclusions}

By mapping the boundary layer dynamics into the orthogonal coordinates $(\eta_1, \eta_2, \eta_3)$, this analysis establishes the following definitive conclusions based on the campaign-specific footprint:

\begin{itemize}
    \item \textbf{ Delayed Threshold with Weak-Transversality Crossing: } The continuation analysis identifies a Hopf instability at a critical dissipation parameter nearly an order of magnitude below the corresponding CASES-99 threshold (gamma\_c approx 0.0235 versus 0.278). The crossing rate magnitude is substantially smaller than the grassland value, indicating a weakly emergent instability rather than an abrupt high-transversality transition. Both observations follow directly from the eigenanalysis.
    \item \textbf{ Dynamically Distinct Oscillatory Timescale: } At the stability boundary the dominant eigenpair crosses the imaginary axis with a non-zero imaginary component corresponding to a characteristic period of approximately 31.5 minutes. This timescale falls within the range commonly associated with wave-mediated and shear-intermittency processes observed in strongly stable boundary layers. The period is substantially shorter than the oscillatory mode identified in CASES-99, suggesting that the instability emerging over snow is dynamically distinct rather than a delayed manifestation of the same mode.
    \item \textbf{ Evidence for Manifold Reshaping: } The combination of a substantially reduced critical threshold and an altered oscillation frequency implies that snow-covered stable boundary layers occupy a fundamentally different region of low-dimensional state space. The threshold and frequency ratios between campaigns are inconsistent with a simple parameter rescaling of the grassland dynamics. These results suggest that surface conditions influence not only the location of stability boundaries but also the geometric structure of the underlying dynamical manifold reconstructed from observations.
\end{itemize}

\subsection{Operational Framework Recommendation}

These signatures confirm that \texttt{ FLOSS } contributes campaign-specific constraints in the verification pipeline. Its operational role can be summarized through runtime geometry and conditioning diagnostics:
\begin{equation}
    \mathcal{M}_{\text{validation}} = \mathbf{f}\Big(\mathcal{A}_{\text{site}}\big(\kappa_{\text{rank}}=48.95,\ D_{\text{eff}}=1.5369\big),\ \mathcal{S}_{\text{surface}}\Big)
\end{equation}

Surface roughness context for this run: \texttt{ 0.001 m }.

When integrated into multi-campaign workflows, low-roughness and canopy-dominated sites jointly provide empirical constraints for model behavior in non-local, intermittent, and collapse-prone stable boundary-layer conditions.

\label{sec:conclusions_FLOSS}
